% sddswput: write waveform data from SDDS files to EPICS PVs.
\begin{sddsprog}{sddswput}
\item \textbf{description:}
  \verb+sddswput+ reads an SDDS file containing waveform data and writes the
  values to the corresponding EPICS process variables.  The input file is
  typically produced by \progref{sddswget} and must contain the parameter
  \verb|WaveformPV| and the column \verb|Waveform|.

\item \textbf{examples:}
  To restore a waveform saved in \verb|kicker.wf|:
  \begin{verbatim}
  sddswput kicker.wf
  
\end{verbatim}

\item \textbf{synopsis:}
  \begin{verbatim}
usage: sddswput <inputfile>
      [-pendIOtime=<seconds>] [-pipe[=input]]
  
\end{verbatim}

\item \textbf{files:}
\begin{itemize}
  \item \textbf{input file:} SDDS file with parameter \verb|WaveformPV| and column \verb|Waveform|.
    An SDDS6 enum \verb|Waveform| column is restored using its dictionary
    labels, rather than its codes. Labels must fit in the Channel Access
    string limit of 39 bytes. Plain string elements longer than 39 bytes are
    truncated to that limit.
\end{itemize}

\item \textbf{switches:}
\begin{itemize}
  \item {\tt -pendIOtime=<seconds>} --- time to wait for Channel Access searches and puts. The default is 10 seconds.
  \item {\tt -pipe[=input]} --- read the input SDDS file from standard input.
\end{itemize}

\item \textbf{see also:}
\begin{itemize}
  \item \progref{sddswmonitor}
\end{itemize}

\item \textbf{author:} Michael Borland, ANL/APS.
\end{sddsprog}
